\begin{longtable}{|
    >{\raggedright\arraybackslash}p{4.3cm} |
    >{\raggedright\arraybackslash}p{4.3cm} |
    >{\raggedright\arraybackslash}p{4.1cm} |}

\caption{Summary of Literature Review}
\label{tab:literature-review} \\
\hline
\multicolumn{1}{|c|}{\parbox[c][2em][c]{4cm}{\centering \textbf{Title}}} &
\multicolumn{1}{c|}{\parbox[c][2em][c]{4.5cm}{\centering \textbf{Methodology}}} &
\multicolumn{1}{c|}{\parbox[c][2em][c]{4.5cm}{\centering \textbf{Key Inferences}}} \\
\hline
\textit{A Survey on Speech Deepfake Detection} (2024) \cite{Li2024_survey} 
& Systematic analysis of over 200 publications (up to March 2024) covering the entire detection pipeline; Evaluates both fully synthetic and partially manipulated audio
& Shift from handcrafted features toward deep embeddings from SSL models like wav2vec 2.0; OC-Softmax as a superior loss function 
\\ \hline
\textit{Towards the Development of a Real-Time Deepfake Audio Detection System in Communication Platforms}
\cite{Mathew2024}
&
Evaluates models trained on ASVspoof 2019 under live Microsoft Teams communication streams. Assesses deployment performance using precision, recall, and F-score metrics.
&
Models trained on offline datasets degrade significantly in real communication environments due to compression and platform artifacts.
\\ \hline
\textit{RTCFake: Speech Deepfake Detection in Real-Time Communication}
\cite{Xue2025_RTCFake}
&
Constructs a 600-hour dataset through real-world black-box transmission across seven communication platforms including Zoom and WeChat. Proposes \textbf{Phoneme-Guided Consistency Learning (PCL)} to learn platform-invariant speech representations.
&
RTC transmission introduces severe domain mismatch not captured by traditional datasets such as ASVspoof. Phoneme-level representations remain more stable than frame-level acoustic features under platform-induced distortions, improving generalization across unseen communication platforms.
\\ \hline
\textit{ASVspoof 5: Crowdsourced Speech Data, Deepfakes, and Adversarial Attacks at Scale}
\cite{Wang2024_asv5}
&
Introduces crowdsourced speech recordings, neural audio codecs, and adversarial attacks targeting both deepfake detection and speaker verification systems.
&
Neural codecs and adversarial perturbations significantly reduce detector reliability; robust calibration is required.
\\ \hline
\textit{Real-Time Detection of AI-Generated Speech for Deepfake Voice Conversion}
\cite{Bird2023}
&
Introduces the DEEP-VOICE dataset and evaluates lightweight machine learning algorithms including Random Forest and XGBoost for short-window real-time inference.
&
XGBoost achieves high detection performance with extremely low inference latency, demonstrating that optimized traditional machine learning approaches remain viable for edge and real-time deployment scenarios.
\\ \hline
\textit{Benchmarking Audio Deepfake Detection Robustness in Real-World Communication Scenarios}
\cite{Shi2024}
&
Introduces the \textbf{ADD-C benchmark}, simulating six communication codecs and packet loss rates ranging from 0\% to 20\%. 
&
Models such as LCNN and AASIST experience performance decreases with increasing packet loss; codec-aware augmentation improves robustness.
\\ \hline
\textit{MultiAPI Spoof: A Multi-API Dataset and Local-Attention Network for Speech Anti-Spoofing Detection}
\cite{Zhang2025_MultiAPI}
&
Constructs a dataset containing spoofed speech generated from 30 different APIs and proposes a Local-Attention enhanced Nes2Net architecture for fine-grained feature learning.
&
Source diversity significantly affects detector generalization. Local-attention mechanisms improve robustness across different synthesis systems, while API tracing introduces attribution as an emerging research direction.
\\ \hline
\end{longtable}
